\section{Discussion: empirical scope and future directions}
\label{sec:discussion}

\Cref{thm:main_convergence_biased_sgd,thm:lower_bound,cor:entropy_decay,cor:pl_exponential_rate}
characterise test-time scaling for any pair $(\text{policy}, F)$ for
which the assumptions of \Cref{sec:assumptions} hold. The formalism is
silent about \emph{which} pairs those are. This section discusses the
empirical scope, falsifiable predictions, honest caveats, and a short
list of optimization-style follow-ups.

\subsection{Empirical scope}
\label{sec:discussion_scope}

The central assumption (gradient correctness with bounded bias,
\Cref{ass:gradient_correctness_bounded_bias}) is the substantive
empirical content of the framework: it pins down a single condition
under which the dynamics of $x_j$ are a biased-SGD recursion on the
explicit loss $\loss(\cdot; Q)$. The remaining four assumptions
(\Cref{ass:bounded_second_moment,ass:bounded_value_norms,ass:unembedding_norm,ass:initial_loss})
are mild structural conditions that hold for any bounded-weight
transformer.

\paragraph{Trained reasoning models satisfy (GC-bb); untrained ones do not.}
The numerical sanity check
(\texttt{risk-2-sanity-check.md} in the project's
\texttt{.proof-research/}) finds that for randomly-initialised
one-layer/one-head transformers, approximately $50\%$ of
$(j, \text{init}, \text{question})$ triples have $\hat \eta_j < 0$ ---
that is, $\E[g_j \mid \mathcal F_{j-1}]$ is on average no better
correlated with $-\nabla \loss$ than with $+\nabla \loss$. There is
\emph{no} $\eta_0 > 0$ for which \Cref{ass:gradient_correctness_bounded_bias}
holds in the untrained regime. A small amount of stylised training
alignment ($\alpha \ge 0.1$, where $\alpha$ is the fraction of the value
vector aligned with the correct-token unembedding direction)
\emph{flips this completely}: the negative-$\hat\eta_j$ fraction drops
to zero and the bias-floor estimate $\beta^2 / \eta_0$ shrinks by a
factor of $\sim 10\times$. The honest characterisation of the
framework's empirical scope is therefore:
\Cref{ass:gradient_correctness_bounded_bias} is a property of
\emph{reasoning-tuned} policies trained on verifiable rewards
\cite{openai2024o1, deepseek2025r1, qwen2025thinking}; it is empirically
falsifiable, but it does \emph{not} hold for untrained transformers and
should be expected to fail on non-reasoning-tuned models.

\paragraph{Tasks the framework targets.}
We designed the framework with verifiable-reward reasoning in mind:
mathematics, code synthesis, and discrete logical-deduction problems
of the kind targeted by
\cite{wei2022cot, openai2024o1, deepseek2025r1, qwen2025thinking}. For
these tasks the correct-token set $\Cset(Q)$ has a natural realisation:
tokens inside the final answer (digits, the contents of a
\verb|\boxed{}|, tokens following an ``\texttt{Answer:}'' marker), or
synonyms thereof. The constrained-softmax loss $\loss(x; Q)$ then
quantifies how concentrated the model's posterior is on
\emph{any} correct decoding, and \Cref{thm:main_convergence_biased_sgd}
shows that biased-SGD-style descent on $\loss$ converges at rate
$O(1/\sqrt T) + O(\beta^2 / \eta_0)$.

\paragraph{Tasks the framework does not cover.}
Two regimes break the assumptions even before any constants are
estimated. \emph{Open-ended generation} (creative writing,
summarisation, free-form dialogue) has no discrete $\Cset(Q)$ to
condition on; the constrained-softmax loss has no natural
instantiation, and the framework as stated does not apply. \emph{Long-horizon
planning and multi-step agentic tasks} localise progress not at the
token level but across tool calls, observations, and subgoal-bearing
tokens, whose value-projection signature need not align with the final
answer's $\nabla \loss$; the framework can in principle be applied
per-subgoal under a hierarchical decomposition, but we do not formalise
that extension.

\subsection{Falsifiable predictions}
\label{sec:discussion_predictions}

The framework yields three concrete, falsifiable predictions about
test-time scaling that can be tested without retraining.

\paragraph{(i) Loss-vs-$T$ shape and the floor-dominated regime.}
\Cref{thm:main_convergence_biased_sgd} predicts that the randomised-iterate
loss $\E[\loss(x_R; Q)]$ at horizon $T$ satisfies
\begin{align*}
   \E[\loss(x_R; Q)]
   \;\asymp\;
   \frac{c_1\, L_0}{\sqrt T}
   \;+\; \frac{c_2\, \beta^2}{\eta_0},
\end{align*}
a sum of a rate term decaying as $1/\sqrt T$ and a $T$-independent
\emph{bias floor} $c_2 \beta^2 / \eta_0$. The numerical sanity check
(\texttt{risk-2-sanity-check.md}) finds that for trained reasoning models
the rate term \emph{dominates} the floor at all empirically-realistic
$T \le 10^5$: the break-even $T$ where the two terms agree is
$\gtrsim 10^9$ in the median trained regime. This is a sharp empirical
prediction: practitioners should expect the loss curve to look like a
$1/\sqrt T$ decay across the entire deployed horizon, with no visible
floor at $T \le 10^5$. Equivalently, the regime where the
biased-SGD floor binds is \emph{below} the typical deployment scale,
making the floor a long-term ($T \gg 10^6$) consideration rather than a
near-term constraint.

\paragraph{(ii) Entropy plateau.}
\Cref{cor:entropy_decay} predicts that the next-token entropy at the
$\langle/\mathrm{think}\rangle$ position decays in expectation at the
square-root of the loss rate
($\E[|H_T - H_\infty(Q)|] \asymp \sqrt{\E[\loss(x_T; Q)]}$ under the
local strong-convexity condition stated in the corollary), and
plateaus at a model-dependent floor $H_\infty(Q)$ once $\loss(x_T; Q)$
hits its bias floor. This is the empirical signal documented by
\cite{choi2025entropy}: the plateau they use as an early-exit trigger is
the empirical signature of $\loss(x_T; Q)$ approaching its bias floor.
The framework predicts both the qualitative shape (decay-then-plateau)
and a quantitative relationship: the plateau level $H_\infty(Q)$ is
determined by the trained-policy parameters $(\beta, \eta_0)$ together
with the unembedding $W_U$.

\paragraph{(iii) Per-policy heterogeneity of $(\eta_0, \beta)$.}
Different reasoning-tuned policies have different values of $\eta_0$
(the lower bound on the per-step alignment of $\E[g_j]$ with $-\nabla \loss$)
and $\beta$ (the supremum of the bias scale $\norm{b_j}/\eta_j$). The
framework predicts that more-aggressive verifiable-reward training
should yield \emph{larger} $\eta_0$ (better alignment) and \emph{smaller}
$\beta$ (smaller residual bias), giving a smaller bias floor and a
faster effective rate. A direct evaluation protocol --- estimating
$(\eta_0, \beta)$ from a few hundred forward passes per policy by the
Monte-Carlo residual procedure of \Cref{rem:gc_bb_remark} --- would
provide a per-model two-number fingerprint of inference-time gradient
quality, comparable across families and training regimes in the way that
loss curves and gradient norms are comparable across training runs.
This is plausibly more diagnostic than pass@1 alone, since pass@1
conflates the rate $\eta_0$ and the floor $\beta^2 / \eta_0$.

\subsection{Honest scope caveats}
\label{sec:discussion_caveats}

We close the empirical-scope discussion with four honest caveats whose
acknowledgement is essential to the framework's sellability.

\paragraph{(a) $\loss$ is non-convex; PL is an extra hypothesis.}
\Cref{rem:smoothness_non_convex} establishes that $\loss(\cdot; Q)$ is
$L_{\mathrm{sm}}$-smooth but \emph{not} convex --- its Hessian subtracts
the $\qcond$-covariance from the $p$-covariance, and the difference can
be negative definite. The PL inequality of \Cref{cor:pl_exponential_rate}
is therefore an \emph{additional} hypothesis on the trajectory's
interaction with the landscape: it cannot be derived from architecture or
from \Cref{ass:gradient_correctness_bounded_bias}.

\paragraph{(b) (GC-bb) is a postulate, not an architectural fact.}
\Cref{rem:gc_bb_remark} establishes that any alignment between
$\E[g_j \mid \mathcal F_{j-1}]$ and $-\nabla \loss$ is learned, not
structural --- there is no architectural identity that forces the value
vector $V_j$ to point in the residual direction $W_U^\top (p - \qcond)$.
The framework therefore diagnoses the policy's gradient quality but
does not derive it from first principles.

\paragraph{(c) $\{\loss \le \log 2\}$ is sufficient, not exact.}
\Cref{rem:sufficient_not_exact} establishes that the level set
$\{\loss \le \log 2\}$ is a smooth interior subset of the polyhedral
decoding-correctness cone $\mathcal D(Q)$, but does not coincide with
it: the cone extends far outside the level set. The framework therefore
gives a sufficient condition for correct decoding, not a tight
characterisation; a future cone-based analysis could potentially
sharpen the failure-probability bound.

\paragraph{(d) Empirically relevant regime is rate-dominated, not floor-dominated.}
The numerical sanity check (\texttt{risk-2-sanity-check.md}, \S 5.2)
finds that at $T \le 10^5$ in the trained-policy median regime
($\alpha = 0.1$ alignment, median $\hat\eta_0 \approx 5 \times 10^{-3}$,
$L_0 \approx 6$), the rate term $c_1 L_0/\sqrt T$ in
\Cref{eq:main_loss_bound} is $\sim 5.5 \times \log 2$ — the bound is
technically non-vacuous, but the rate term remains above the
Markov-via-logit-margin threshold of $\log 2$, so the
failure-probability bound \Cref{eq:main_failure_bound} does not yield a
non-trivial empirical prediction at the deployment horizon. (One would
need $T \gtrsim 3 \times 10^6$ for the rate term to drop below $\log 2$
in this regime.) This is why the PL strengthening
(\Cref{cor:pl_exponential_rate}) is the empirically-relevant statement:
the exponential rate is what test-time scaling curves of
\cite{openai2024o1, deepseek2025r1} qualitatively resemble, and the
framework's headline statement at the empirical horizon is therefore
best read as \emph{conditional} on the PL hypothesis. The honest interpretation: the main theorem
(\Cref{thm:main_convergence_biased_sgd}) sets up the framework; the PL
corollary (\Cref{cor:pl_exponential_rate}) provides the
practitioner-relevant rate; the open empirical question of whether PL
holds on inference trajectories is what gates the framework's predictive
power at deployment scale.

\subsection{Future directions}
\label{sec:discussion_future}

We close with four directions that take the optimization-shaped framing
seriously, in the spirit of ``this is the next paper, not this one''.

\paragraph{Optimization-style diagnostics from $(\eta_0, \beta)$.}
The constants $(\eta_0, \beta)$ are estimable from any open-weights
reasoning model via the residual procedure of \Cref{rem:gc_bb_remark}.
Reporting $(\eta_0, \beta)$ alongside pass@1 gives a structural
breakdown of inference-time gradient quality, comparable across
families in the way that gradient-noise scale and effective batch size
are comparable across training runs. We expect this two-number
fingerprint to be more diagnostic than pass@1 for studying
inference-time-compute trade-offs, since pass@1 conflates the rate
$\eta_0$ and the floor $\beta^2/\eta_0$ into a single performance scalar.

\paragraph{Momentum-style acceleration at inference time.}
The recurrence $x_j = (s_{j-1}/s_j) x_{j-1} + g_j$ of
\Cref{lem:softmax_running_average} has the algebraic form of an SGD
step with effective weight-decay $1 - s_{j-1}/s_j$ on an explicit loss
$\loss(\cdot; Q)$. Standard convex-optimisation acceleration --- Polyak /
Nesterov momentum, adaptive step-size schedules, Adam-style second-moment
scaling --- operates by reshaping the effective step-size schedule.
Applied at inference time as a re-weighting of the attention softmax
(e.g., a temperature schedule that reshapes $s_j$ as a function of $j$),
can these techniques produce a strictly faster within-inference rate?
The open issue is whether the reshaping preserves
\Cref{ass:gradient_correctness_bounded_bias} or breaks it.

\paragraph{Compounding within-chain and across-sample scaling.}
\Cref{thm:main_convergence_biased_sgd} governs failure probability
\emph{within} a single reasoning chain of length $T$. A complementary
line of work considers parallel sampling: draw $N$ independent chains
and aggregate by majority vote or best-of-$N$ verifier scoring; the
failure probability then decays in $N$. Combining the two regimes, a
practitioner with a total compute budget of $N T$ forward passes can
trade $T$ (longer thinking) against $N$ (more samples) and pick the
operating point that minimises the joint failure rate. A clean compound
bound formalising this trade-off --- and predicting the optimal $(N, T)$
allocation as a function of $(\eta_0, \beta)$ --- would be a directly
useful extension for practitioners deploying reasoning models at fixed
inference budgets.

\paragraph{Empirically validating the PL inequality on inference trajectories.}
\Cref{cor:pl_exponential_rate} is the empirically-relevant statement at
deployment-scale $T$, but its hypothesis (PL on the trajectory's basin)
is an \emph{open empirical question}. A direct test would estimate the
PL constant $\mu$ along a sample of inference trajectories (by fitting
$\log(\loss(x_j) - \hat L_*)$ versus $j$ for an estimated minimum loss
$\hat L_*$), and check whether the predicted exponential rate
$\exp(-c_5 \eta_0 \mu T)$ matches the observed pass@1 curve at the same
horizon. Such an experiment, applied across a fixed model suite (e.g.\
DeepSeek-R1, o1, Qwen3-thinking, plus a non-reasoning-tuned control),
would test the framework's empirical predictive power and identify
the policies for which PL is a useful description versus those for
which it fails.
